% ======================================================================
%  AT-MONO111 — Chapter 12: Cosmology
%  Canonical Monograph (v2.0) · The Actualization Theory:
%  A Reconstruction of Physics from Difference, Actualization and Spectrum
%
%  Status: COMPLETE — PASS-READY (MONO007 referee-ready architecture)
%  Inputs: Chapter01_Difference.tex ... Chapter11_StandardModel.tex
%  Method: assembly only — canonical results only, no new physics, no speculation
%  Citations: QG227 (initial conditions), QG228 (information content),
%             QG230 (lambda origin), QG231 (structure formation),
%             QG234 (density fractions), QG237 (CMB spectrum),
%             QG238 (acoustic peaks), QG240 (blind reproduction),
%             QG295 (spectrum necessity), QG296 (reconstruction)
%
%  Usage: % ======================================================================
%  AT-MONO111 — Chapter 12: Cosmology
%  Canonical Monograph (v2.0) · The Actualization Theory:
%  A Reconstruction of Physics from Difference, Actualization and Spectrum
%
%  Status: COMPLETE — PASS-READY (MONO007 referee-ready architecture)
%  Inputs: Chapter01_Difference.tex ... Chapter11_StandardModel.tex
%  Method: assembly only — canonical results only, no new physics, no speculation
%  Citations: QG227 (initial conditions), QG228 (information content),
%             QG230 (lambda origin), QG231 (structure formation),
%             QG234 (density fractions), QG237 (CMB spectrum),
%             QG238 (acoustic peaks), QG240 (blind reproduction),
%             QG295 (spectrum necessity), QG296 (reconstruction)
%
%  Usage: % ======================================================================
%  AT-MONO111 — Chapter 12: Cosmology
%  Canonical Monograph (v2.0) · The Actualization Theory:
%  A Reconstruction of Physics from Difference, Actualization and Spectrum
%
%  Status: COMPLETE — PASS-READY (MONO007 referee-ready architecture)
%  Inputs: Chapter01_Difference.tex ... Chapter11_StandardModel.tex
%  Method: assembly only — canonical results only, no new physics, no speculation
%  Citations: QG227 (initial conditions), QG228 (information content),
%             QG230 (lambda origin), QG231 (structure formation),
%             QG234 (density fractions), QG237 (CMB spectrum),
%             QG238 (acoustic peaks), QG240 (blind reproduction),
%             QG295 (spectrum necessity), QG296 (reconstruction)
%
%  Usage: % ======================================================================
%  AT-MONO111 — Chapter 12: Cosmology
%  Canonical Monograph (v2.0) · The Actualization Theory:
%  A Reconstruction of Physics from Difference, Actualization and Spectrum
%
%  Status: COMPLETE — PASS-READY (MONO007 referee-ready architecture)
%  Inputs: Chapter01_Difference.tex ... Chapter11_StandardModel.tex
%  Method: assembly only — canonical results only, no new physics, no speculation
%  Citations: QG227 (initial conditions), QG228 (information content),
%             QG230 (lambda origin), QG231 (structure formation),
%             QG234 (density fractions), QG237 (CMB spectrum),
%             QG238 (acoustic peaks), QG240 (blind reproduction),
%             QG295 (spectrum necessity), QG296 (reconstruction)
%
%  Usage: \input{Chapter12_Cosmology.tex} after the preamble
%         that defines the theorem environments (or include via the master file).
% ======================================================================

% ---- Theorem environments (define in the master preamble if not present) ----
% \newtheorem{definition}{Definition}[chapter]
% \newtheorem{lemma}[definition]{Lemma}
% \newtheorem{theorem}[definition]{Theorem}
% \newtheorem{corollary}[definition]{Corollary}
% \newtheorem{remark}[definition]{Remark}

\chapter{Cosmology}
\label{chap:cosmology}

\section{Introduction}
\label{sec:cosmology-introduction}

The preceding chapters derived the quantum and gravitational structure and the standard
model of The Actualization Theory from the canonical foundation. This chapter completes the
physics part of the monograph by deriving \emph{cosmology} from the same foundation: the
cosmological constant, the matter and vacuum fractions, the cosmic microwave background, the
structure formation, and the expansion history. Every cosmological observable of this
chapter is traced back to the D96 spectral structure \cite{QG230,QG234,QG237,QG238}.

The central results are: the cosmological constant $\Lambda$ is the residual actualization
pressure of the critical branching vacuum \cite{QG230}; the density fractions
$\Omega_\Lambda$ and $\Omega_m$ are the information-density fractions of the D96 octave
record \cite{QG234}; the CMB acoustic structure is the standing-wave harmonic structure of
the D96 recombination-scale mode ladder \cite{QG237,QG238}; structure formation follows
the actualization dynamics and the spectral hierarchy \cite{QG231}; and the expansion
history is a spectrum consequence, requiring no independent cosmological primitive
\cite{QG227,QG228,QG240}.

Throughout, cosmology is presented as \emph{emergent} --- a spectrum readout --- and no
claims beyond the accepted derivations are made. In particular, no inflation claims beyond
the accepted results, no new cosmological mechanisms, no unresolved Bekenstein discussion,
and no boundary material are included (the boundary topics are reserved for the Boundary
Layer part of the monograph). We use only accepted canonical results and introduce no new
physics and no speculation.

\section{Cosmology as a Spectrum Readout}
\label{sec:spectrum-readout}

\begin{definition}[Cosmological readout]
\label{def:cosmo-readout}
A \emph{cosmological readout} is a derived observable of the theory obtained as a function
of the D96 spectral structure: the cosmological quantities are reads of the spectrum, in
the same sense as the physics sectors of the preceding chapters.
\end{definition}

\begin{theorem}[Cosmological observables emerge from the D96 spectral structure]
\label{thm:cosmo-emergent}
The cosmological observables of the theory emerge from the D96 spectral structure: they are
derived functions of the spectral content, not independent inputs.
\end{theorem}

\begin{proof}
The reconstruction audit \cite{QG296} reconstructs the major results of the theory ---
including the cosmological ones --- from the minimal hierarchy, with the spectrum layer as
the source of the constants used by all sectors. The cosmological quantities are reads of
the spectral content: $\Lambda$ from the branching-vacuum structure \cite{QG230}, the
fractions from the octave record \cite{QG234}, the CMB structure from the mode ladder
\cite{QG237,QG238}. Hence cosmology is a spectrum readout --- emergent, not primitive.
\end{proof}

\begin{corollary}[Cosmology is emergent]
\label{cor:cosmo-emergent}
Cosmology is emergent in the theory: its observables are reads of the D96 spectrum, which
is itself the derived output of the actualization attractor.
\end{corollary}

\begin{proof}
By Theorem~\ref{thm:cosmo-emergent}, the cosmological observables are functions of the
spectral structure; by Chapter~\ref{chap:d96} (Theorem~\ref{thm:d96-derived}), the spectrum
is the derived output of the attractor. Hence cosmology is twice-emergent, consistent with
the canonical architecture.
\end{proof}

\section{The Cosmological Constant}
\label{sec:lambda}

\begin{definition}[Cosmological constant]
\label{def:lambda}
The \emph{cosmological constant} $\Lambda$ is the residual actualization pressure of the
critical branching vacuum: the realized vacuum's departure from its uniform expectation.
\end{definition}

\begin{theorem}[Existence, sign, and scaling of \texorpdfstring{$\Lambda$}{Lambda}]
\label{thm:lambda}
The cosmological constant exists, is positive, and scales with the spectral structure: at
criticality ($\mu=1$) the Galton--Watson mean is constant while the variance grows
($\mathrm{Var}(Z_k) = k\sigma^2$), the realized vacuum never equals its uniform expectation,
and the residual pressure $\Lambda$ follows.
\end{theorem}

\begin{proof}
The lambda-origin audit \cite{QG230} establishes the construction. Existence: at criticality
the Galton--Watson mean is constant but the variance grows, giving a residual pressure.
Sign: the realized vacuum carries positive information relative to the uniform expectation
(the information-content result \cite{QG228}), so the residual pressure is positive.
Scaling: the residual pressure is a function of the branching-vacuum structure. Hence
$\Lambda$ exists, is positive, and scales with the actualization structure --- derived, not
imported.
\end{proof}

\begin{corollary}[\texorpdfstring{$\Lambda$}{Lambda} is emergent, not fundamental]
\label{cor:lambda-emergent}
The cosmological constant is emergent, not fundamental: it is the residual pressure of the
actualization vacuum, not an independent parameter.
\end{corollary}

\begin{proof}
By Theorem~\ref{thm:lambda}, $\Lambda$ is the residual pressure of the critical branching
vacuum. By Theorem~\ref{thm:cosmo-emergent}, it is a read of the spectral/actualization
structure. Hence $\Lambda$ is emergent --- a derived consequence of the actualization
process, not a fundamental input.
\end{proof}

\section{Matter and Vacuum Fractions}
\label{sec:fractions}

\begin{definition}[Density fractions]
\label{def:fractions}
The \emph{density fractions} $\Omega_\Lambda$ and $\Omega_m$ are the information-density
fractions of the D96 octave record: the partition of the total density between the vacuum
and the matter content.
\end{definition}

\begin{theorem}[The density fractions are derived from spectral occupancies]
\label{thm:fractions}
The cosmological density fractions are derived from the D96 octave occupancies:
$\Omega_\Lambda = I_{\mathrm{occ}}/\ln K$ and $\Omega_m = 1 - \Omega_\Lambda$, where
$I_{\mathrm{occ}}$ is the realized octave record's information content.
\end{theorem}

\begin{proof}
The fraction-origin audit \cite{QG234} establishes the construction: the fractions are the
information-density fractions of the D96 octave record, with
$\Omega_\Lambda = I_{\mathrm{occ}}/\ln K = 0.6839$ and $\Omega_m = 0.3161$, matching
observation within a fraction of a percent. The information content
$I_{\mathrm{occ}} = KL(p\Vert\mathrm{uniform}) = 0.7513$ nats is computed from the D96
occupancies $[4,4,87]$ with $95$ modes \cite{QG228}. Hence the fractions are derived from
the spectral occupancies.
\end{proof}

\begin{corollary}[The density partition follows the spectrum structure]
\label{cor:density-partition}
The density partition of the universe follows the spectrum structure: the vacuum fraction
is the octave-record information, and the matter fraction its complement.
\end{corollary}

\begin{proof}
By Theorem~\ref{thm:fractions}, $\Omega_\Lambda$ is the octave-record information fraction
and $\Omega_m$ its complement. Both are functions of the D96 occupancies
(Chapter~\ref{chap:d96}, Corollary~\ref{cor:constants-derived}). Hence the density
partition follows the spectrum structure.
\end{proof}

\section{The Cosmic Microwave Background}
\label{sec:cmb}

\begin{definition}[CMB acoustic structure]
\label{def:cmb}
The \emph{CMB acoustic structure} is the standing-wave harmonic structure of the
recombination-scale field: the acoustic peaks are the harmonics of the D96 mode ladder.
\end{definition}

\begin{theorem}[The acoustic structure is derived from the spectrum]
\label{thm:acoustic}
The CMB acoustic structure is the standing-wave harmonic structure of the D96
recombination-scale mode ladder: the acoustic peaks are the harmonics of the spectral
mode structure.
\end{theorem}

\begin{proof}
The acoustic-peak-origin audit \cite{QG238} establishes the construction: the acoustic peaks
are the standing-wave harmonics of the recombination-scale field, which is the D96 octave
spectrum's mode ladder. The first peak $\ell_1 = 220.48$ and the ratios
$r_{21} = 2.4368$, $r_{31} = 3.6965$ are derived from the D96 octave hierarchy. Hence the
CMB acoustic structure is derived from the spectrum.
\end{proof}

\begin{corollary}[The peak hierarchy follows spectral organization]
\label{cor:peak-hierarchy}
The acoustic-peak hierarchy follows the spectral organization: the peak positions and
ratios are the harmonics and ratios of the D96 mode ladder.
\end{corollary}

\begin{proof}
By Theorem~\ref{thm:acoustic}, the peaks are the harmonics of the D96 mode ladder; the peak
ratios are ratios of the spectral harmonics
(Chapter~\ref{chap:d96}, Corollary~\ref{cor:constants-derived}). Hence the peak hierarchy
follows the spectral organization.
\end{proof}

\begin{lemma}[The seed spectrum is the counting variance]
\label{lem:seed-spectrum}
The seed power spectrum is the Poisson counting variance $\delta_i = 1/\sqrt{\langle N
\rangle}$, scale-free from critical branching.
\end{lemma}

\begin{proof}
The CMB-origin audit \cite{QG237} establishes the seed: the scalar spectral index
$n_s$ is the octave-hierarchy tilt of the D96 spectrum, and the seed power spectrum is the
Poisson counting variance, scale-free ($n_s = 1$) from critical branching. Hence the seed
spectrum is the counting variance of the actualization process.
\end{proof}

\section{Structure Formation}
\label{sec:structure}

\begin{theorem}[Large-scale organization follows actualization dynamics]
\label{thm:structure-formation}
The large-scale organization of the universe follows the actualization dynamics: structure
formation is derived from the Q-event statistics and the spectral hierarchy.
\end{theorem}

\begin{proof}
The structure-formation audit \cite{QG231} establishes the construction: structure
formation is derived from the Q-event statistics of the actualization process, with no new
primitives. The initial conditions are the uniform critical state \cite{QG227}, and the
information content provides the departure that seeds the structure \cite{QG228}. Hence
large-scale organization follows the actualization dynamics.
\end{proof}

\begin{lemma}[Matter clustering is consistent with the spectral hierarchy]
\label{lem:clustering}
The matter clustering of the universe is consistent with the spectral hierarchy: the
clustering scale structure is a read of the D96 spectral organization.
\end{lemma}

\begin{proof}
The structure-formation result \cite{QG231} and the CMB-acoustic result \cite{QG238}
together establish the consistency: the seed spectrum is the counting variance
(Lemma~\ref{lem:seed-spectrum}), the acoustic structure is the spectral harmonics
(Theorem~\ref{thm:acoustic}), and the clustering follows the same spectral hierarchy. Hence
matter clustering is consistent with the spectral organization.
\end{proof}

\section{The Expansion History}
\label{sec:expansion}

\begin{theorem}[Cosmological evolution is a spectrum consequence]
\label{thm:expansion}
The cosmological expansion is a spectrum consequence: the evolution of the universe is a
read of the actualization structure, requiring no independent cosmological primitive.
\end{theorem}

\begin{proof}
The cosmological observables --- the constant $\Lambda$ (Theorem~\ref{thm:lambda}), the
fractions (Theorem~\ref{thm:fractions}), the CMB structure (Theorem~\ref{thm:acoustic}),
and the structure formation (Theorem~\ref{thm:structure-formation}) --- are all reads of
the spectral/actualization structure. The blind reproduction audit \cite{QG240} confirms
that the cosmological quantities are recomputed from the D96 quantities alone, with no
fitted inputs. Hence the expansion history is a spectrum consequence.
\end{proof}

\begin{corollary}[No independent cosmological primitive]
\label{cor:no-cosmo-primitive}
Cosmology introduces no independent primitive: every cosmological observable is a read of
the D96 spectrum and the actualization process.
\end{corollary}

\begin{proof}
By Theorem~\ref{thm:expansion}, the cosmological evolution is a spectrum consequence; by
Theorem~\ref{thm:cosmo-emergent}, the observables are functions of the spectral content.
Hence cosmology adds no primitive to the canonical foundation
$\{\text{Difference},\eta\}$.
\end{proof}

\section{Consequences}
\label{sec:cosmo-consequences}

\begin{theorem}[Cosmology is consistent with the canonical architecture]
\label{thm:cosmo-consistent}
Cosmology is consistent with the canonical architecture: every cosmological observable is
a read of the D96 spectral moments, with the physics part of the monograph now complete.
\end{theorem}

\begin{proof}
The cosmological constant (Theorem~\ref{thm:lambda}), the fractions
(Theorem~\ref{thm:fractions}), the CMB structure (Theorem~\ref{thm:acoustic}), the structure
formation (Theorem~\ref{thm:structure-formation}), and the expansion
(Theorem~\ref{thm:expansion}) are all reads of the D96 spectral/actualization structure,
consistent with the spectrum-necessity result \cite{QG295} and the reconstruction audit
\cite{QG296}. Hence cosmology is consistent with the canonical architecture, completing the
physics part.
\end{proof}

\begin{corollary}[Difference → Actualization → Spectrum → Cosmology]
\label{cor:canonical-cosmology}
The cosmological sector completes the canonical chain:
\[
\text{Difference}\;\longrightarrow\;\text{Actualization}\;\longrightarrow\;\text{Inevitable Spectrum}\;\longrightarrow\;\text{Cosmology},
\]
with every cosmological observable a read of the spectrum.
\end{corollary}

\begin{proof}
By Theorem~\ref{thm:cosmo-consistent}, the cosmological observables are reads of the D96
spectrum; by Chapter~\ref{chap:spectrum} (Theorem~\ref{thm:core-completed}), the spectrum
is the output of the actualization process rooted in Difference. Hence the cosmological
sector is the terminal read of the canonical chain.
\end{proof}

\begin{remark}[Referee-safe wording]
Consistent with MONO007 \cite{MONO007}, this chapter presents cosmology as emergent and makes
no claims beyond the accepted derivations. No inflation claims beyond the accepted results
(the initial-condition and information origins replace the inflation epoch), no new
cosmological mechanisms, no unresolved Bekenstein discussion, and no boundary material are
included; the boundary topics are reserved for the Boundary Layer part of the monograph.
\end{remark}

\section{Summary}
\label{sec:cosmo-summary}

This chapter has derived cosmology. We have proved:

\begin{enumerate}
\item \textbf{Cosmology as a spectrum readout} (Theorem~\ref{thm:cosmo-emergent},
Corollary~\ref{cor:cosmo-emergent}): the cosmological observables emerge from the D96
spectral structure, and cosmology is emergent;
\item \textbf{The cosmological constant} (Theorem~\ref{thm:lambda},
Corollary~\ref{cor:lambda-emergent}): $\Lambda$ exists, is positive, and scales with the
actualization structure --- emergent, not fundamental;
\item \textbf{Matter and vacuum fractions} (Theorem~\ref{thm:fractions},
Corollary~\ref{cor:density-partition}): $\Omega_\Lambda$ and $\Omega_m$ are derived from
the D96 octave occupancies;
\item \textbf{The cosmic microwave background} (Theorem~\ref{thm:acoustic},
Corollary~\ref{cor:peak-hierarchy}, Lemma~\ref{lem:seed-spectrum}): the acoustic structure
is the D96 mode-ladder harmonics, with the seed spectrum the counting variance;
\item \textbf{Structure formation} (Theorem~\ref{thm:structure-formation},
Lemma~\ref{lem:clustering}): large-scale organization follows the actualization dynamics,
consistent with the spectral hierarchy;
\item \textbf{The expansion history} (Theorem~\ref{thm:expansion},
Corollary~\ref{cor:no-cosmo-primitive}): cosmological evolution is a spectrum consequence,
with no independent cosmological primitive;
\item \textbf{Consequences} (Theorem~\ref{thm:cosmo-consistent},
Corollary~\ref{cor:canonical-cosmology}): cosmology is consistent with the canonical
architecture, completing the physics part of the monograph.
\end{enumerate}

Cosmology is therefore emergent in The Actualization Theory: the cosmological constant is
the residual actualization pressure, the density fractions are the octave-record
information, the CMB acoustic structure is the spectral harmonics, structure formation
follows the actualization dynamics, and the expansion history is a spectrum consequence ---
with every observable traced back to the D96 spectral moments and no independent
cosmological primitive. This completes the physics part of the monograph
(Difference → Actualization → Spectrum → Physics): the foundation, the spectrum, quantum
mechanics, gravity, the standard model, and cosmology are all derived reads of the one
canonical structure.

\begin{thebibliography}{99}
\bibitem{QG227} AT-QG Phase 227, \emph{Initial Conditions Origin}, Docs/Research.
\bibitem{QG228} AT-QG Phase 228, \emph{Information Content Origin}, Docs/Research.
\bibitem{QG230} AT-QG Phase 230, \emph{Lambda Origin}, Docs/Research.
\bibitem{QG231} AT-QG Phase 231, \emph{Structure Formation Origin}, Docs/Research.
\bibitem{QG234} AT-QG Phase 234, \emph{Cosmological Fractions Origin}, Docs/Research.
\bibitem{QG237} AT-QG Phase 237, \emph{CMB Spectrum Origin}, Docs/Research.
\bibitem{QG238} AT-QG Phase 238, \emph{Acoustic Peak Origin}, Docs/Research.
\bibitem{QG240} AT-QG Phase 240, \emph{Cosmology Blind Reproduction}, Docs/Research.
\bibitem{QG295} AT-QG Phase 295, \emph{Spectrum Necessity Audit}, Docs/Research.
\bibitem{QG296} AT-QG Phase 296, \emph{Reconstruction Audit}, Docs/Research.
\bibitem{MONO007} AT-MONO007, \emph{Referee-Readiness Resolution}, Docs/Zenodo/Monograph.
\end{thebibliography}
 after the preamble
%         that defines the theorem environments (or include via the master file).
% ======================================================================

% ---- Theorem environments (define in the master preamble if not present) ----
% \newtheorem{definition}{Definition}[chapter]
% \newtheorem{lemma}[definition]{Lemma}
% \newtheorem{theorem}[definition]{Theorem}
% \newtheorem{corollary}[definition]{Corollary}
% \newtheorem{remark}[definition]{Remark}

\chapter{Cosmology}
\label{chap:cosmology}

\section{Introduction}
\label{sec:cosmology-introduction}

The preceding chapters derived the quantum and gravitational structure and the standard
model of The Actualization Theory from the canonical foundation. This chapter completes the
physics part of the monograph by deriving \emph{cosmology} from the same foundation: the
cosmological constant, the matter and vacuum fractions, the cosmic microwave background, the
structure formation, and the expansion history. Every cosmological observable of this
chapter is traced back to the D96 spectral structure \cite{QG230,QG234,QG237,QG238}.

The central results are: the cosmological constant $\Lambda$ is the residual actualization
pressure of the critical branching vacuum \cite{QG230}; the density fractions
$\Omega_\Lambda$ and $\Omega_m$ are the information-density fractions of the D96 octave
record \cite{QG234}; the CMB acoustic structure is the standing-wave harmonic structure of
the D96 recombination-scale mode ladder \cite{QG237,QG238}; structure formation follows
the actualization dynamics and the spectral hierarchy \cite{QG231}; and the expansion
history is a spectrum consequence, requiring no independent cosmological primitive
\cite{QG227,QG228,QG240}.

Throughout, cosmology is presented as \emph{emergent} --- a spectrum readout --- and no
claims beyond the accepted derivations are made. In particular, no inflation claims beyond
the accepted results, no new cosmological mechanisms, no unresolved Bekenstein discussion,
and no boundary material are included (the boundary topics are reserved for the Boundary
Layer part of the monograph). We use only accepted canonical results and introduce no new
physics and no speculation.

\section{Cosmology as a Spectrum Readout}
\label{sec:spectrum-readout}

\begin{definition}[Cosmological readout]
\label{def:cosmo-readout}
A \emph{cosmological readout} is a derived observable of the theory obtained as a function
of the D96 spectral structure: the cosmological quantities are reads of the spectrum, in
the same sense as the physics sectors of the preceding chapters.
\end{definition}

\begin{theorem}[Cosmological observables emerge from the D96 spectral structure]
\label{thm:cosmo-emergent}
The cosmological observables of the theory emerge from the D96 spectral structure: they are
derived functions of the spectral content, not independent inputs.
\end{theorem}

\begin{proof}
The reconstruction audit \cite{QG296} reconstructs the major results of the theory ---
including the cosmological ones --- from the minimal hierarchy, with the spectrum layer as
the source of the constants used by all sectors. The cosmological quantities are reads of
the spectral content: $\Lambda$ from the branching-vacuum structure \cite{QG230}, the
fractions from the octave record \cite{QG234}, the CMB structure from the mode ladder
\cite{QG237,QG238}. Hence cosmology is a spectrum readout --- emergent, not primitive.
\end{proof}

\begin{corollary}[Cosmology is emergent]
\label{cor:cosmo-emergent}
Cosmology is emergent in the theory: its observables are reads of the D96 spectrum, which
is itself the derived output of the actualization attractor.
\end{corollary}

\begin{proof}
By Theorem~\ref{thm:cosmo-emergent}, the cosmological observables are functions of the
spectral structure; by Chapter~\ref{chap:d96} (Theorem~\ref{thm:d96-derived}), the spectrum
is the derived output of the attractor. Hence cosmology is twice-emergent, consistent with
the canonical architecture.
\end{proof}

\section{The Cosmological Constant}
\label{sec:lambda}

\begin{definition}[Cosmological constant]
\label{def:lambda}
The \emph{cosmological constant} $\Lambda$ is the residual actualization pressure of the
critical branching vacuum: the realized vacuum's departure from its uniform expectation.
\end{definition}

\begin{theorem}[Existence, sign, and scaling of \texorpdfstring{$\Lambda$}{Lambda}]
\label{thm:lambda}
The cosmological constant exists, is positive, and scales with the spectral structure: at
criticality ($\mu=1$) the Galton--Watson mean is constant while the variance grows
($\mathrm{Var}(Z_k) = k\sigma^2$), the realized vacuum never equals its uniform expectation,
and the residual pressure $\Lambda$ follows.
\end{theorem}

\begin{proof}
The lambda-origin audit \cite{QG230} establishes the construction. Existence: at criticality
the Galton--Watson mean is constant but the variance grows, giving a residual pressure.
Sign: the realized vacuum carries positive information relative to the uniform expectation
(the information-content result \cite{QG228}), so the residual pressure is positive.
Scaling: the residual pressure is a function of the branching-vacuum structure. Hence
$\Lambda$ exists, is positive, and scales with the actualization structure --- derived, not
imported.
\end{proof}

\begin{corollary}[\texorpdfstring{$\Lambda$}{Lambda} is emergent, not fundamental]
\label{cor:lambda-emergent}
The cosmological constant is emergent, not fundamental: it is the residual pressure of the
actualization vacuum, not an independent parameter.
\end{corollary}

\begin{proof}
By Theorem~\ref{thm:lambda}, $\Lambda$ is the residual pressure of the critical branching
vacuum. By Theorem~\ref{thm:cosmo-emergent}, it is a read of the spectral/actualization
structure. Hence $\Lambda$ is emergent --- a derived consequence of the actualization
process, not a fundamental input.
\end{proof}

\section{Matter and Vacuum Fractions}
\label{sec:fractions}

\begin{definition}[Density fractions]
\label{def:fractions}
The \emph{density fractions} $\Omega_\Lambda$ and $\Omega_m$ are the information-density
fractions of the D96 octave record: the partition of the total density between the vacuum
and the matter content.
\end{definition}

\begin{theorem}[The density fractions are derived from spectral occupancies]
\label{thm:fractions}
The cosmological density fractions are derived from the D96 octave occupancies:
$\Omega_\Lambda = I_{\mathrm{occ}}/\ln K$ and $\Omega_m = 1 - \Omega_\Lambda$, where
$I_{\mathrm{occ}}$ is the realized octave record's information content.
\end{theorem}

\begin{proof}
The fraction-origin audit \cite{QG234} establishes the construction: the fractions are the
information-density fractions of the D96 octave record, with
$\Omega_\Lambda = I_{\mathrm{occ}}/\ln K = 0.6839$ and $\Omega_m = 0.3161$, matching
observation within a fraction of a percent. The information content
$I_{\mathrm{occ}} = KL(p\Vert\mathrm{uniform}) = 0.7513$ nats is computed from the D96
occupancies $[4,4,87]$ with $95$ modes \cite{QG228}. Hence the fractions are derived from
the spectral occupancies.
\end{proof}

\begin{corollary}[The density partition follows the spectrum structure]
\label{cor:density-partition}
The density partition of the universe follows the spectrum structure: the vacuum fraction
is the octave-record information, and the matter fraction its complement.
\end{corollary}

\begin{proof}
By Theorem~\ref{thm:fractions}, $\Omega_\Lambda$ is the octave-record information fraction
and $\Omega_m$ its complement. Both are functions of the D96 occupancies
(Chapter~\ref{chap:d96}, Corollary~\ref{cor:constants-derived}). Hence the density
partition follows the spectrum structure.
\end{proof}

\section{The Cosmic Microwave Background}
\label{sec:cmb}

\begin{definition}[CMB acoustic structure]
\label{def:cmb}
The \emph{CMB acoustic structure} is the standing-wave harmonic structure of the
recombination-scale field: the acoustic peaks are the harmonics of the D96 mode ladder.
\end{definition}

\begin{theorem}[The acoustic structure is derived from the spectrum]
\label{thm:acoustic}
The CMB acoustic structure is the standing-wave harmonic structure of the D96
recombination-scale mode ladder: the acoustic peaks are the harmonics of the spectral
mode structure.
\end{theorem}

\begin{proof}
The acoustic-peak-origin audit \cite{QG238} establishes the construction: the acoustic peaks
are the standing-wave harmonics of the recombination-scale field, which is the D96 octave
spectrum's mode ladder. The first peak $\ell_1 = 220.48$ and the ratios
$r_{21} = 2.4368$, $r_{31} = 3.6965$ are derived from the D96 octave hierarchy. Hence the
CMB acoustic structure is derived from the spectrum.
\end{proof}

\begin{corollary}[The peak hierarchy follows spectral organization]
\label{cor:peak-hierarchy}
The acoustic-peak hierarchy follows the spectral organization: the peak positions and
ratios are the harmonics and ratios of the D96 mode ladder.
\end{corollary}

\begin{proof}
By Theorem~\ref{thm:acoustic}, the peaks are the harmonics of the D96 mode ladder; the peak
ratios are ratios of the spectral harmonics
(Chapter~\ref{chap:d96}, Corollary~\ref{cor:constants-derived}). Hence the peak hierarchy
follows the spectral organization.
\end{proof}

\begin{lemma}[The seed spectrum is the counting variance]
\label{lem:seed-spectrum}
The seed power spectrum is the Poisson counting variance $\delta_i = 1/\sqrt{\langle N
\rangle}$, scale-free from critical branching.
\end{lemma}

\begin{proof}
The CMB-origin audit \cite{QG237} establishes the seed: the scalar spectral index
$n_s$ is the octave-hierarchy tilt of the D96 spectrum, and the seed power spectrum is the
Poisson counting variance, scale-free ($n_s = 1$) from critical branching. Hence the seed
spectrum is the counting variance of the actualization process.
\end{proof}

\section{Structure Formation}
\label{sec:structure}

\begin{theorem}[Large-scale organization follows actualization dynamics]
\label{thm:structure-formation}
The large-scale organization of the universe follows the actualization dynamics: structure
formation is derived from the Q-event statistics and the spectral hierarchy.
\end{theorem}

\begin{proof}
The structure-formation audit \cite{QG231} establishes the construction: structure
formation is derived from the Q-event statistics of the actualization process, with no new
primitives. The initial conditions are the uniform critical state \cite{QG227}, and the
information content provides the departure that seeds the structure \cite{QG228}. Hence
large-scale organization follows the actualization dynamics.
\end{proof}

\begin{lemma}[Matter clustering is consistent with the spectral hierarchy]
\label{lem:clustering}
The matter clustering of the universe is consistent with the spectral hierarchy: the
clustering scale structure is a read of the D96 spectral organization.
\end{lemma}

\begin{proof}
The structure-formation result \cite{QG231} and the CMB-acoustic result \cite{QG238}
together establish the consistency: the seed spectrum is the counting variance
(Lemma~\ref{lem:seed-spectrum}), the acoustic structure is the spectral harmonics
(Theorem~\ref{thm:acoustic}), and the clustering follows the same spectral hierarchy. Hence
matter clustering is consistent with the spectral organization.
\end{proof}

\section{The Expansion History}
\label{sec:expansion}

\begin{theorem}[Cosmological evolution is a spectrum consequence]
\label{thm:expansion}
The cosmological expansion is a spectrum consequence: the evolution of the universe is a
read of the actualization structure, requiring no independent cosmological primitive.
\end{theorem}

\begin{proof}
The cosmological observables --- the constant $\Lambda$ (Theorem~\ref{thm:lambda}), the
fractions (Theorem~\ref{thm:fractions}), the CMB structure (Theorem~\ref{thm:acoustic}),
and the structure formation (Theorem~\ref{thm:structure-formation}) --- are all reads of
the spectral/actualization structure. The blind reproduction audit \cite{QG240} confirms
that the cosmological quantities are recomputed from the D96 quantities alone, with no
fitted inputs. Hence the expansion history is a spectrum consequence.
\end{proof}

\begin{corollary}[No independent cosmological primitive]
\label{cor:no-cosmo-primitive}
Cosmology introduces no independent primitive: every cosmological observable is a read of
the D96 spectrum and the actualization process.
\end{corollary}

\begin{proof}
By Theorem~\ref{thm:expansion}, the cosmological evolution is a spectrum consequence; by
Theorem~\ref{thm:cosmo-emergent}, the observables are functions of the spectral content.
Hence cosmology adds no primitive to the canonical foundation
$\{\text{Difference},\eta\}$.
\end{proof}

\section{Consequences}
\label{sec:cosmo-consequences}

\begin{theorem}[Cosmology is consistent with the canonical architecture]
\label{thm:cosmo-consistent}
Cosmology is consistent with the canonical architecture: every cosmological observable is
a read of the D96 spectral moments, with the physics part of the monograph now complete.
\end{theorem}

\begin{proof}
The cosmological constant (Theorem~\ref{thm:lambda}), the fractions
(Theorem~\ref{thm:fractions}), the CMB structure (Theorem~\ref{thm:acoustic}), the structure
formation (Theorem~\ref{thm:structure-formation}), and the expansion
(Theorem~\ref{thm:expansion}) are all reads of the D96 spectral/actualization structure,
consistent with the spectrum-necessity result \cite{QG295} and the reconstruction audit
\cite{QG296}. Hence cosmology is consistent with the canonical architecture, completing the
physics part.
\end{proof}

\begin{corollary}[Difference → Actualization → Spectrum → Cosmology]
\label{cor:canonical-cosmology}
The cosmological sector completes the canonical chain:
\[
\text{Difference}\;\longrightarrow\;\text{Actualization}\;\longrightarrow\;\text{Inevitable Spectrum}\;\longrightarrow\;\text{Cosmology},
\]
with every cosmological observable a read of the spectrum.
\end{corollary}

\begin{proof}
By Theorem~\ref{thm:cosmo-consistent}, the cosmological observables are reads of the D96
spectrum; by Chapter~\ref{chap:spectrum} (Theorem~\ref{thm:core-completed}), the spectrum
is the output of the actualization process rooted in Difference. Hence the cosmological
sector is the terminal read of the canonical chain.
\end{proof}

\begin{remark}[Referee-safe wording]
Consistent with MONO007 \cite{MONO007}, this chapter presents cosmology as emergent and makes
no claims beyond the accepted derivations. No inflation claims beyond the accepted results
(the initial-condition and information origins replace the inflation epoch), no new
cosmological mechanisms, no unresolved Bekenstein discussion, and no boundary material are
included; the boundary topics are reserved for the Boundary Layer part of the monograph.
\end{remark}

\section{Summary}
\label{sec:cosmo-summary}

This chapter has derived cosmology. We have proved:

\begin{enumerate}
\item \textbf{Cosmology as a spectrum readout} (Theorem~\ref{thm:cosmo-emergent},
Corollary~\ref{cor:cosmo-emergent}): the cosmological observables emerge from the D96
spectral structure, and cosmology is emergent;
\item \textbf{The cosmological constant} (Theorem~\ref{thm:lambda},
Corollary~\ref{cor:lambda-emergent}): $\Lambda$ exists, is positive, and scales with the
actualization structure --- emergent, not fundamental;
\item \textbf{Matter and vacuum fractions} (Theorem~\ref{thm:fractions},
Corollary~\ref{cor:density-partition}): $\Omega_\Lambda$ and $\Omega_m$ are derived from
the D96 octave occupancies;
\item \textbf{The cosmic microwave background} (Theorem~\ref{thm:acoustic},
Corollary~\ref{cor:peak-hierarchy}, Lemma~\ref{lem:seed-spectrum}): the acoustic structure
is the D96 mode-ladder harmonics, with the seed spectrum the counting variance;
\item \textbf{Structure formation} (Theorem~\ref{thm:structure-formation},
Lemma~\ref{lem:clustering}): large-scale organization follows the actualization dynamics,
consistent with the spectral hierarchy;
\item \textbf{The expansion history} (Theorem~\ref{thm:expansion},
Corollary~\ref{cor:no-cosmo-primitive}): cosmological evolution is a spectrum consequence,
with no independent cosmological primitive;
\item \textbf{Consequences} (Theorem~\ref{thm:cosmo-consistent},
Corollary~\ref{cor:canonical-cosmology}): cosmology is consistent with the canonical
architecture, completing the physics part of the monograph.
\end{enumerate}

Cosmology is therefore emergent in The Actualization Theory: the cosmological constant is
the residual actualization pressure, the density fractions are the octave-record
information, the CMB acoustic structure is the spectral harmonics, structure formation
follows the actualization dynamics, and the expansion history is a spectrum consequence ---
with every observable traced back to the D96 spectral moments and no independent
cosmological primitive. This completes the physics part of the monograph
(Difference → Actualization → Spectrum → Physics): the foundation, the spectrum, quantum
mechanics, gravity, the standard model, and cosmology are all derived reads of the one
canonical structure.

\begin{thebibliography}{99}
\bibitem{QG227} AT-QG Phase 227, \emph{Initial Conditions Origin}, Docs/Research.
\bibitem{QG228} AT-QG Phase 228, \emph{Information Content Origin}, Docs/Research.
\bibitem{QG230} AT-QG Phase 230, \emph{Lambda Origin}, Docs/Research.
\bibitem{QG231} AT-QG Phase 231, \emph{Structure Formation Origin}, Docs/Research.
\bibitem{QG234} AT-QG Phase 234, \emph{Cosmological Fractions Origin}, Docs/Research.
\bibitem{QG237} AT-QG Phase 237, \emph{CMB Spectrum Origin}, Docs/Research.
\bibitem{QG238} AT-QG Phase 238, \emph{Acoustic Peak Origin}, Docs/Research.
\bibitem{QG240} AT-QG Phase 240, \emph{Cosmology Blind Reproduction}, Docs/Research.
\bibitem{QG295} AT-QG Phase 295, \emph{Spectrum Necessity Audit}, Docs/Research.
\bibitem{QG296} AT-QG Phase 296, \emph{Reconstruction Audit}, Docs/Research.
\bibitem{MONO007} AT-MONO007, \emph{Referee-Readiness Resolution}, Docs/Zenodo/Monograph.
\end{thebibliography}
 after the preamble
%         that defines the theorem environments (or include via the master file).
% ======================================================================

% ---- Theorem environments (define in the master preamble if not present) ----
% \newtheorem{definition}{Definition}[chapter]
% \newtheorem{lemma}[definition]{Lemma}
% \newtheorem{theorem}[definition]{Theorem}
% \newtheorem{corollary}[definition]{Corollary}
% \newtheorem{remark}[definition]{Remark}

\chapter{Cosmology}
\label{chap:cosmology}

\section{Introduction}
\label{sec:cosmology-introduction}

The preceding chapters derived the quantum and gravitational structure and the standard
model of The Actualization Theory from the canonical foundation. This chapter completes the
physics part of the monograph by deriving \emph{cosmology} from the same foundation: the
cosmological constant, the matter and vacuum fractions, the cosmic microwave background, the
structure formation, and the expansion history. Every cosmological observable of this
chapter is traced back to the D96 spectral structure \cite{QG230,QG234,QG237,QG238}.

The central results are: the cosmological constant $\Lambda$ is the residual actualization
pressure of the critical branching vacuum \cite{QG230}; the density fractions
$\Omega_\Lambda$ and $\Omega_m$ are the information-density fractions of the D96 octave
record \cite{QG234}; the CMB acoustic structure is the standing-wave harmonic structure of
the D96 recombination-scale mode ladder \cite{QG237,QG238}; structure formation follows
the actualization dynamics and the spectral hierarchy \cite{QG231}; and the expansion
history is a spectrum consequence, requiring no independent cosmological primitive
\cite{QG227,QG228,QG240}.

Throughout, cosmology is presented as \emph{emergent} --- a spectrum readout --- and no
claims beyond the accepted derivations are made. In particular, no inflation claims beyond
the accepted results, no new cosmological mechanisms, no unresolved Bekenstein discussion,
and no boundary material are included (the boundary topics are reserved for the Boundary
Layer part of the monograph). We use only accepted canonical results and introduce no new
physics and no speculation.

\section{Cosmology as a Spectrum Readout}
\label{sec:spectrum-readout}

\begin{definition}[Cosmological readout]
\label{def:cosmo-readout}
A \emph{cosmological readout} is a derived observable of the theory obtained as a function
of the D96 spectral structure: the cosmological quantities are reads of the spectrum, in
the same sense as the physics sectors of the preceding chapters.
\end{definition}

\begin{theorem}[Cosmological observables emerge from the D96 spectral structure]
\label{thm:cosmo-emergent}
The cosmological observables of the theory emerge from the D96 spectral structure: they are
derived functions of the spectral content, not independent inputs.
\end{theorem}

\begin{proof}
The reconstruction audit \cite{QG296} reconstructs the major results of the theory ---
including the cosmological ones --- from the minimal hierarchy, with the spectrum layer as
the source of the constants used by all sectors. The cosmological quantities are reads of
the spectral content: $\Lambda$ from the branching-vacuum structure \cite{QG230}, the
fractions from the octave record \cite{QG234}, the CMB structure from the mode ladder
\cite{QG237,QG238}. Hence cosmology is a spectrum readout --- emergent, not primitive.
\end{proof}

\begin{corollary}[Cosmology is emergent]
\label{cor:cosmo-emergent}
Cosmology is emergent in the theory: its observables are reads of the D96 spectrum, which
is itself the derived output of the actualization attractor.
\end{corollary}

\begin{proof}
By Theorem~\ref{thm:cosmo-emergent}, the cosmological observables are functions of the
spectral structure; by Chapter~\ref{chap:d96} (Theorem~\ref{thm:d96-derived}), the spectrum
is the derived output of the attractor. Hence cosmology is twice-emergent, consistent with
the canonical architecture.
\end{proof}

\section{The Cosmological Constant}
\label{sec:lambda}

\begin{definition}[Cosmological constant]
\label{def:lambda}
The \emph{cosmological constant} $\Lambda$ is the residual actualization pressure of the
critical branching vacuum: the realized vacuum's departure from its uniform expectation.
\end{definition}

\begin{theorem}[Existence, sign, and scaling of \texorpdfstring{$\Lambda$}{Lambda}]
\label{thm:lambda}
The cosmological constant exists, is positive, and scales with the spectral structure: at
criticality ($\mu=1$) the Galton--Watson mean is constant while the variance grows
($\mathrm{Var}(Z_k) = k\sigma^2$), the realized vacuum never equals its uniform expectation,
and the residual pressure $\Lambda$ follows.
\end{theorem}

\begin{proof}
The lambda-origin audit \cite{QG230} establishes the construction. Existence: at criticality
the Galton--Watson mean is constant but the variance grows, giving a residual pressure.
Sign: the realized vacuum carries positive information relative to the uniform expectation
(the information-content result \cite{QG228}), so the residual pressure is positive.
Scaling: the residual pressure is a function of the branching-vacuum structure. Hence
$\Lambda$ exists, is positive, and scales with the actualization structure --- derived, not
imported.
\end{proof}

\begin{corollary}[\texorpdfstring{$\Lambda$}{Lambda} is emergent, not fundamental]
\label{cor:lambda-emergent}
The cosmological constant is emergent, not fundamental: it is the residual pressure of the
actualization vacuum, not an independent parameter.
\end{corollary}

\begin{proof}
By Theorem~\ref{thm:lambda}, $\Lambda$ is the residual pressure of the critical branching
vacuum. By Theorem~\ref{thm:cosmo-emergent}, it is a read of the spectral/actualization
structure. Hence $\Lambda$ is emergent --- a derived consequence of the actualization
process, not a fundamental input.
\end{proof}

\section{Matter and Vacuum Fractions}
\label{sec:fractions}

\begin{definition}[Density fractions]
\label{def:fractions}
The \emph{density fractions} $\Omega_\Lambda$ and $\Omega_m$ are the information-density
fractions of the D96 octave record: the partition of the total density between the vacuum
and the matter content.
\end{definition}

\begin{theorem}[The density fractions are derived from spectral occupancies]
\label{thm:fractions}
The cosmological density fractions are derived from the D96 octave occupancies:
$\Omega_\Lambda = I_{\mathrm{occ}}/\ln K$ and $\Omega_m = 1 - \Omega_\Lambda$, where
$I_{\mathrm{occ}}$ is the realized octave record's information content.
\end{theorem}

\begin{proof}
The fraction-origin audit \cite{QG234} establishes the construction: the fractions are the
information-density fractions of the D96 octave record, with
$\Omega_\Lambda = I_{\mathrm{occ}}/\ln K = 0.6839$ and $\Omega_m = 0.3161$, matching
observation within a fraction of a percent. The information content
$I_{\mathrm{occ}} = KL(p\Vert\mathrm{uniform}) = 0.7513$ nats is computed from the D96
occupancies $[4,4,87]$ with $95$ modes \cite{QG228}. Hence the fractions are derived from
the spectral occupancies.
\end{proof}

\begin{corollary}[The density partition follows the spectrum structure]
\label{cor:density-partition}
The density partition of the universe follows the spectrum structure: the vacuum fraction
is the octave-record information, and the matter fraction its complement.
\end{corollary}

\begin{proof}
By Theorem~\ref{thm:fractions}, $\Omega_\Lambda$ is the octave-record information fraction
and $\Omega_m$ its complement. Both are functions of the D96 occupancies
(Chapter~\ref{chap:d96}, Corollary~\ref{cor:constants-derived}). Hence the density
partition follows the spectrum structure.
\end{proof}

\section{The Cosmic Microwave Background}
\label{sec:cmb}

\begin{definition}[CMB acoustic structure]
\label{def:cmb}
The \emph{CMB acoustic structure} is the standing-wave harmonic structure of the
recombination-scale field: the acoustic peaks are the harmonics of the D96 mode ladder.
\end{definition}

\begin{theorem}[The acoustic structure is derived from the spectrum]
\label{thm:acoustic}
The CMB acoustic structure is the standing-wave harmonic structure of the D96
recombination-scale mode ladder: the acoustic peaks are the harmonics of the spectral
mode structure.
\end{theorem}

\begin{proof}
The acoustic-peak-origin audit \cite{QG238} establishes the construction: the acoustic peaks
are the standing-wave harmonics of the recombination-scale field, which is the D96 octave
spectrum's mode ladder. The first peak $\ell_1 = 220.48$ and the ratios
$r_{21} = 2.4368$, $r_{31} = 3.6965$ are derived from the D96 octave hierarchy. Hence the
CMB acoustic structure is derived from the spectrum.
\end{proof}

\begin{corollary}[The peak hierarchy follows spectral organization]
\label{cor:peak-hierarchy}
The acoustic-peak hierarchy follows the spectral organization: the peak positions and
ratios are the harmonics and ratios of the D96 mode ladder.
\end{corollary}

\begin{proof}
By Theorem~\ref{thm:acoustic}, the peaks are the harmonics of the D96 mode ladder; the peak
ratios are ratios of the spectral harmonics
(Chapter~\ref{chap:d96}, Corollary~\ref{cor:constants-derived}). Hence the peak hierarchy
follows the spectral organization.
\end{proof}

\begin{lemma}[The seed spectrum is the counting variance]
\label{lem:seed-spectrum}
The seed power spectrum is the Poisson counting variance $\delta_i = 1/\sqrt{\langle N
\rangle}$, scale-free from critical branching.
\end{lemma}

\begin{proof}
The CMB-origin audit \cite{QG237} establishes the seed: the scalar spectral index
$n_s$ is the octave-hierarchy tilt of the D96 spectrum, and the seed power spectrum is the
Poisson counting variance, scale-free ($n_s = 1$) from critical branching. Hence the seed
spectrum is the counting variance of the actualization process.
\end{proof}

\section{Structure Formation}
\label{sec:structure}

\begin{theorem}[Large-scale organization follows actualization dynamics]
\label{thm:structure-formation}
The large-scale organization of the universe follows the actualization dynamics: structure
formation is derived from the Q-event statistics and the spectral hierarchy.
\end{theorem}

\begin{proof}
The structure-formation audit \cite{QG231} establishes the construction: structure
formation is derived from the Q-event statistics of the actualization process, with no new
primitives. The initial conditions are the uniform critical state \cite{QG227}, and the
information content provides the departure that seeds the structure \cite{QG228}. Hence
large-scale organization follows the actualization dynamics.
\end{proof}

\begin{lemma}[Matter clustering is consistent with the spectral hierarchy]
\label{lem:clustering}
The matter clustering of the universe is consistent with the spectral hierarchy: the
clustering scale structure is a read of the D96 spectral organization.
\end{lemma}

\begin{proof}
The structure-formation result \cite{QG231} and the CMB-acoustic result \cite{QG238}
together establish the consistency: the seed spectrum is the counting variance
(Lemma~\ref{lem:seed-spectrum}), the acoustic structure is the spectral harmonics
(Theorem~\ref{thm:acoustic}), and the clustering follows the same spectral hierarchy. Hence
matter clustering is consistent with the spectral organization.
\end{proof}

\section{The Expansion History}
\label{sec:expansion}

\begin{theorem}[Cosmological evolution is a spectrum consequence]
\label{thm:expansion}
The cosmological expansion is a spectrum consequence: the evolution of the universe is a
read of the actualization structure, requiring no independent cosmological primitive.
\end{theorem}

\begin{proof}
The cosmological observables --- the constant $\Lambda$ (Theorem~\ref{thm:lambda}), the
fractions (Theorem~\ref{thm:fractions}), the CMB structure (Theorem~\ref{thm:acoustic}),
and the structure formation (Theorem~\ref{thm:structure-formation}) --- are all reads of
the spectral/actualization structure. The blind reproduction audit \cite{QG240} confirms
that the cosmological quantities are recomputed from the D96 quantities alone, with no
fitted inputs. Hence the expansion history is a spectrum consequence.
\end{proof}

\begin{corollary}[No independent cosmological primitive]
\label{cor:no-cosmo-primitive}
Cosmology introduces no independent primitive: every cosmological observable is a read of
the D96 spectrum and the actualization process.
\end{corollary}

\begin{proof}
By Theorem~\ref{thm:expansion}, the cosmological evolution is a spectrum consequence; by
Theorem~\ref{thm:cosmo-emergent}, the observables are functions of the spectral content.
Hence cosmology adds no primitive to the canonical foundation
$\{\text{Difference},\eta\}$.
\end{proof}

\section{Consequences}
\label{sec:cosmo-consequences}

\begin{theorem}[Cosmology is consistent with the canonical architecture]
\label{thm:cosmo-consistent}
Cosmology is consistent with the canonical architecture: every cosmological observable is
a read of the D96 spectral moments, with the physics part of the monograph now complete.
\end{theorem}

\begin{proof}
The cosmological constant (Theorem~\ref{thm:lambda}), the fractions
(Theorem~\ref{thm:fractions}), the CMB structure (Theorem~\ref{thm:acoustic}), the structure
formation (Theorem~\ref{thm:structure-formation}), and the expansion
(Theorem~\ref{thm:expansion}) are all reads of the D96 spectral/actualization structure,
consistent with the spectrum-necessity result \cite{QG295} and the reconstruction audit
\cite{QG296}. Hence cosmology is consistent with the canonical architecture, completing the
physics part.
\end{proof}

\begin{corollary}[Difference → Actualization → Spectrum → Cosmology]
\label{cor:canonical-cosmology}
The cosmological sector completes the canonical chain:
\[
\text{Difference}\;\longrightarrow\;\text{Actualization}\;\longrightarrow\;\text{Inevitable Spectrum}\;\longrightarrow\;\text{Cosmology},
\]
with every cosmological observable a read of the spectrum.
\end{corollary}

\begin{proof}
By Theorem~\ref{thm:cosmo-consistent}, the cosmological observables are reads of the D96
spectrum; by Chapter~\ref{chap:spectrum} (Theorem~\ref{thm:core-completed}), the spectrum
is the output of the actualization process rooted in Difference. Hence the cosmological
sector is the terminal read of the canonical chain.
\end{proof}

\begin{remark}[Referee-safe wording]
Consistent with MONO007 \cite{MONO007}, this chapter presents cosmology as emergent and makes
no claims beyond the accepted derivations. No inflation claims beyond the accepted results
(the initial-condition and information origins replace the inflation epoch), no new
cosmological mechanisms, no unresolved Bekenstein discussion, and no boundary material are
included; the boundary topics are reserved for the Boundary Layer part of the monograph.
\end{remark}

\section{Summary}
\label{sec:cosmo-summary}

This chapter has derived cosmology. We have proved:

\begin{enumerate}
\item \textbf{Cosmology as a spectrum readout} (Theorem~\ref{thm:cosmo-emergent},
Corollary~\ref{cor:cosmo-emergent}): the cosmological observables emerge from the D96
spectral structure, and cosmology is emergent;
\item \textbf{The cosmological constant} (Theorem~\ref{thm:lambda},
Corollary~\ref{cor:lambda-emergent}): $\Lambda$ exists, is positive, and scales with the
actualization structure --- emergent, not fundamental;
\item \textbf{Matter and vacuum fractions} (Theorem~\ref{thm:fractions},
Corollary~\ref{cor:density-partition}): $\Omega_\Lambda$ and $\Omega_m$ are derived from
the D96 octave occupancies;
\item \textbf{The cosmic microwave background} (Theorem~\ref{thm:acoustic},
Corollary~\ref{cor:peak-hierarchy}, Lemma~\ref{lem:seed-spectrum}): the acoustic structure
is the D96 mode-ladder harmonics, with the seed spectrum the counting variance;
\item \textbf{Structure formation} (Theorem~\ref{thm:structure-formation},
Lemma~\ref{lem:clustering}): large-scale organization follows the actualization dynamics,
consistent with the spectral hierarchy;
\item \textbf{The expansion history} (Theorem~\ref{thm:expansion},
Corollary~\ref{cor:no-cosmo-primitive}): cosmological evolution is a spectrum consequence,
with no independent cosmological primitive;
\item \textbf{Consequences} (Theorem~\ref{thm:cosmo-consistent},
Corollary~\ref{cor:canonical-cosmology}): cosmology is consistent with the canonical
architecture, completing the physics part of the monograph.
\end{enumerate}

Cosmology is therefore emergent in The Actualization Theory: the cosmological constant is
the residual actualization pressure, the density fractions are the octave-record
information, the CMB acoustic structure is the spectral harmonics, structure formation
follows the actualization dynamics, and the expansion history is a spectrum consequence ---
with every observable traced back to the D96 spectral moments and no independent
cosmological primitive. This completes the physics part of the monograph
(Difference → Actualization → Spectrum → Physics): the foundation, the spectrum, quantum
mechanics, gravity, the standard model, and cosmology are all derived reads of the one
canonical structure.

\begin{thebibliography}{99}
\bibitem{QG227} AT-QG Phase 227, \emph{Initial Conditions Origin}, Docs/Research.
\bibitem{QG228} AT-QG Phase 228, \emph{Information Content Origin}, Docs/Research.
\bibitem{QG230} AT-QG Phase 230, \emph{Lambda Origin}, Docs/Research.
\bibitem{QG231} AT-QG Phase 231, \emph{Structure Formation Origin}, Docs/Research.
\bibitem{QG234} AT-QG Phase 234, \emph{Cosmological Fractions Origin}, Docs/Research.
\bibitem{QG237} AT-QG Phase 237, \emph{CMB Spectrum Origin}, Docs/Research.
\bibitem{QG238} AT-QG Phase 238, \emph{Acoustic Peak Origin}, Docs/Research.
\bibitem{QG240} AT-QG Phase 240, \emph{Cosmology Blind Reproduction}, Docs/Research.
\bibitem{QG295} AT-QG Phase 295, \emph{Spectrum Necessity Audit}, Docs/Research.
\bibitem{QG296} AT-QG Phase 296, \emph{Reconstruction Audit}, Docs/Research.
\bibitem{MONO007} AT-MONO007, \emph{Referee-Readiness Resolution}, Docs/Zenodo/Monograph.
\end{thebibliography}
 after the preamble
%         that defines the theorem environments (or include via the master file).
% ======================================================================

% ---- Theorem environments (define in the master preamble if not present) ----
% \newtheorem{definition}{Definition}[chapter]
% \newtheorem{lemma}[definition]{Lemma}
% \newtheorem{theorem}[definition]{Theorem}
% \newtheorem{corollary}[definition]{Corollary}
% \newtheorem{remark}[definition]{Remark}

\chapter{Cosmology}
\label{chap:cosmology}

\section{Introduction}
\label{sec:cosmology-introduction}

The preceding chapters derived the quantum and gravitational structure and the standard
model of The Actualization Theory from the canonical foundation. This chapter completes the
physics part of the monograph by deriving \emph{cosmology} from the same foundation: the
cosmological constant, the matter and vacuum fractions, the cosmic microwave background, the
structure formation, and the expansion history. Every cosmological observable of this
chapter is traced back to the D96 spectral structure \cite{QG230,QG234,QG237,QG238}.

The central results are: the cosmological constant $\Lambda$ is the residual actualization
pressure of the critical branching vacuum \cite{QG230}; the density fractions
$\Omega_\Lambda$ and $\Omega_m$ are the information-density fractions of the D96 octave
record \cite{QG234}; the CMB acoustic structure is the standing-wave harmonic structure of
the D96 recombination-scale mode ladder \cite{QG237,QG238}; structure formation follows
the actualization dynamics and the spectral hierarchy \cite{QG231}; and the expansion
history is a spectrum consequence, requiring no independent cosmological primitive
\cite{QG227,QG228,QG240}.

Throughout, cosmology is presented as \emph{emergent} --- a spectrum readout --- and no
claims beyond the accepted derivations are made. In particular, no inflation claims beyond
the accepted results, no new cosmological mechanisms, no unresolved Bekenstein discussion,
and no boundary material are included (the boundary topics are reserved for the Boundary
Layer part of the monograph). We use only accepted canonical results and introduce no new
physics and no speculation.

\section{Cosmology as a Spectrum Readout}
\label{sec:spectrum-readout}

\begin{definition}[Cosmological readout]
\label{def:cosmo-readout}
A \emph{cosmological readout} is a derived observable of the theory obtained as a function
of the D96 spectral structure: the cosmological quantities are reads of the spectrum, in
the same sense as the physics sectors of the preceding chapters.
\end{definition}

\begin{theorem}[Cosmological observables emerge from the D96 spectral structure]
\label{thm:cosmo-emergent}
The cosmological observables of the theory emerge from the D96 spectral structure: they are
derived functions of the spectral content, not independent inputs.
\end{theorem}

\begin{proof}
The reconstruction audit \cite{QG296} reconstructs the major results of the theory ---
including the cosmological ones --- from the minimal hierarchy, with the spectrum layer as
the source of the constants used by all sectors. The cosmological quantities are reads of
the spectral content: $\Lambda$ from the branching-vacuum structure \cite{QG230}, the
fractions from the octave record \cite{QG234}, the CMB structure from the mode ladder
\cite{QG237,QG238}. Hence cosmology is a spectrum readout --- emergent, not primitive.
\end{proof}

\begin{corollary}[Cosmology is emergent]
\label{cor:cosmo-emergent}
Cosmology is emergent in the theory: its observables are reads of the D96 spectrum, which
is itself the derived output of the actualization attractor.
\end{corollary}

\begin{proof}
By Theorem~\ref{thm:cosmo-emergent}, the cosmological observables are functions of the
spectral structure; by Chapter~\ref{chap:d96} (Theorem~\ref{thm:d96-derived}), the spectrum
is the derived output of the attractor. Hence cosmology is twice-emergent, consistent with
the canonical architecture.
\end{proof}

\section{The Cosmological Constant}
\label{sec:lambda}

\begin{definition}[Cosmological constant]
\label{def:lambda}
The \emph{cosmological constant} $\Lambda$ is the residual actualization pressure of the
critical branching vacuum: the realized vacuum's departure from its uniform expectation.
\end{definition}

\begin{theorem}[Existence, sign, and scaling of \texorpdfstring{$\Lambda$}{Lambda}]
\label{thm:lambda}
The cosmological constant exists, is positive, and scales with the spectral structure: at
criticality ($\mu=1$) the Galton--Watson mean is constant while the variance grows
($\mathrm{Var}(Z_k) = k\sigma^2$), the realized vacuum never equals its uniform expectation,
and the residual pressure $\Lambda$ follows.
\end{theorem}

\begin{proof}
The lambda-origin audit \cite{QG230} establishes the construction. Existence: at criticality
the Galton--Watson mean is constant but the variance grows, giving a residual pressure.
Sign: the realized vacuum carries positive information relative to the uniform expectation
(the information-content result \cite{QG228}), so the residual pressure is positive.
Scaling: the residual pressure is a function of the branching-vacuum structure. Hence
$\Lambda$ exists, is positive, and scales with the actualization structure --- derived, not
imported.
\end{proof}

\begin{corollary}[\texorpdfstring{$\Lambda$}{Lambda} is emergent, not fundamental]
\label{cor:lambda-emergent}
The cosmological constant is emergent, not fundamental: it is the residual pressure of the
actualization vacuum, not an independent parameter.
\end{corollary}

\begin{proof}
By Theorem~\ref{thm:lambda}, $\Lambda$ is the residual pressure of the critical branching
vacuum. By Theorem~\ref{thm:cosmo-emergent}, it is a read of the spectral/actualization
structure. Hence $\Lambda$ is emergent --- a derived consequence of the actualization
process, not a fundamental input.
\end{proof}

\section{Matter and Vacuum Fractions}
\label{sec:fractions}

\begin{definition}[Density fractions]
\label{def:fractions}
The \emph{density fractions} $\Omega_\Lambda$ and $\Omega_m$ are the information-density
fractions of the D96 octave record: the partition of the total density between the vacuum
and the matter content.
\end{definition}

\begin{theorem}[The density fractions are derived from spectral occupancies]
\label{thm:fractions}
The cosmological density fractions are derived from the D96 octave occupancies:
$\Omega_\Lambda = I_{\mathrm{occ}}/\ln K$ and $\Omega_m = 1 - \Omega_\Lambda$, where
$I_{\mathrm{occ}}$ is the realized octave record's information content.
\end{theorem}

\begin{proof}
The fraction-origin audit \cite{QG234} establishes the construction: the fractions are the
information-density fractions of the D96 octave record, with
$\Omega_\Lambda = I_{\mathrm{occ}}/\ln K = 0.6839$ and $\Omega_m = 0.3161$, matching
observation within a fraction of a percent. The information content
$I_{\mathrm{occ}} = KL(p\Vert\mathrm{uniform}) = 0.7513$ nats is computed from the D96
occupancies $[4,4,87]$ with $95$ modes \cite{QG228}. Hence the fractions are derived from
the spectral occupancies.
\end{proof}

\begin{corollary}[The density partition follows the spectrum structure]
\label{cor:density-partition}
The density partition of the universe follows the spectrum structure: the vacuum fraction
is the octave-record information, and the matter fraction its complement.
\end{corollary}

\begin{proof}
By Theorem~\ref{thm:fractions}, $\Omega_\Lambda$ is the octave-record information fraction
and $\Omega_m$ its complement. Both are functions of the D96 occupancies
(Chapter~\ref{chap:d96}, Corollary~\ref{cor:constants-derived}). Hence the density
partition follows the spectrum structure.
\end{proof}

\section{The Cosmic Microwave Background}
\label{sec:cmb}

\begin{definition}[CMB acoustic structure]
\label{def:cmb}
The \emph{CMB acoustic structure} is the standing-wave harmonic structure of the
recombination-scale field: the acoustic peaks are the harmonics of the D96 mode ladder.
\end{definition}

\begin{theorem}[The acoustic structure is derived from the spectrum]
\label{thm:acoustic}
The CMB acoustic structure is the standing-wave harmonic structure of the D96
recombination-scale mode ladder: the acoustic peaks are the harmonics of the spectral
mode structure.
\end{theorem}

\begin{proof}
The acoustic-peak-origin audit \cite{QG238} establishes the construction: the acoustic peaks
are the standing-wave harmonics of the recombination-scale field, which is the D96 octave
spectrum's mode ladder. The first peak $\ell_1 = 220.48$ and the ratios
$r_{21} = 2.4368$, $r_{31} = 3.6965$ are derived from the D96 octave hierarchy. Hence the
CMB acoustic structure is derived from the spectrum.
\end{proof}

\begin{corollary}[The peak hierarchy follows spectral organization]
\label{cor:peak-hierarchy}
The acoustic-peak hierarchy follows the spectral organization: the peak positions and
ratios are the harmonics and ratios of the D96 mode ladder.
\end{corollary}

\begin{proof}
By Theorem~\ref{thm:acoustic}, the peaks are the harmonics of the D96 mode ladder; the peak
ratios are ratios of the spectral harmonics
(Chapter~\ref{chap:d96}, Corollary~\ref{cor:constants-derived}). Hence the peak hierarchy
follows the spectral organization.
\end{proof}

\begin{lemma}[The seed spectrum is the counting variance]
\label{lem:seed-spectrum}
The seed power spectrum is the Poisson counting variance $\delta_i = 1/\sqrt{\langle N
\rangle}$, scale-free from critical branching.
\end{lemma}

\begin{proof}
The CMB-origin audit \cite{QG237} establishes the seed: the scalar spectral index
$n_s$ is the octave-hierarchy tilt of the D96 spectrum, and the seed power spectrum is the
Poisson counting variance, scale-free ($n_s = 1$) from critical branching. Hence the seed
spectrum is the counting variance of the actualization process.
\end{proof}

\section{Structure Formation}
\label{sec:structure}

\begin{theorem}[Large-scale organization follows actualization dynamics]
\label{thm:structure-formation}
The large-scale organization of the universe follows the actualization dynamics: structure
formation is derived from the Q-event statistics and the spectral hierarchy.
\end{theorem}

\begin{proof}
The structure-formation audit \cite{QG231} establishes the construction: structure
formation is derived from the Q-event statistics of the actualization process, with no new
primitives. The initial conditions are the uniform critical state \cite{QG227}, and the
information content provides the departure that seeds the structure \cite{QG228}. Hence
large-scale organization follows the actualization dynamics.
\end{proof}

\begin{lemma}[Matter clustering is consistent with the spectral hierarchy]
\label{lem:clustering}
The matter clustering of the universe is consistent with the spectral hierarchy: the
clustering scale structure is a read of the D96 spectral organization.
\end{lemma}

\begin{proof}
The structure-formation result \cite{QG231} and the CMB-acoustic result \cite{QG238}
together establish the consistency: the seed spectrum is the counting variance
(Lemma~\ref{lem:seed-spectrum}), the acoustic structure is the spectral harmonics
(Theorem~\ref{thm:acoustic}), and the clustering follows the same spectral hierarchy. Hence
matter clustering is consistent with the spectral organization.
\end{proof}

\section{The Expansion History}
\label{sec:expansion}

\begin{theorem}[Cosmological evolution is a spectrum consequence]
\label{thm:expansion}
The cosmological expansion is a spectrum consequence: the evolution of the universe is a
read of the actualization structure, requiring no independent cosmological primitive.
\end{theorem}

\begin{proof}
The cosmological observables --- the constant $\Lambda$ (Theorem~\ref{thm:lambda}), the
fractions (Theorem~\ref{thm:fractions}), the CMB structure (Theorem~\ref{thm:acoustic}),
and the structure formation (Theorem~\ref{thm:structure-formation}) --- are all reads of
the spectral/actualization structure. The blind reproduction audit \cite{QG240} confirms
that the cosmological quantities are recomputed from the D96 quantities alone, with no
fitted inputs. Hence the expansion history is a spectrum consequence.
\end{proof}

\begin{corollary}[No independent cosmological primitive]
\label{cor:no-cosmo-primitive}
Cosmology introduces no independent primitive: every cosmological observable is a read of
the D96 spectrum and the actualization process.
\end{corollary}

\begin{proof}
By Theorem~\ref{thm:expansion}, the cosmological evolution is a spectrum consequence; by
Theorem~\ref{thm:cosmo-emergent}, the observables are functions of the spectral content.
Hence cosmology adds no primitive to the canonical foundation
$\{\text{Difference},\eta\}$.
\end{proof}

\section{Consequences}
\label{sec:cosmo-consequences}

\begin{theorem}[Cosmology is consistent with the canonical architecture]
\label{thm:cosmo-consistent}
Cosmology is consistent with the canonical architecture: every cosmological observable is
a read of the D96 spectral moments, with the physics part of the monograph now complete.
\end{theorem}

\begin{proof}
The cosmological constant (Theorem~\ref{thm:lambda}), the fractions
(Theorem~\ref{thm:fractions}), the CMB structure (Theorem~\ref{thm:acoustic}), the structure
formation (Theorem~\ref{thm:structure-formation}), and the expansion
(Theorem~\ref{thm:expansion}) are all reads of the D96 spectral/actualization structure,
consistent with the spectrum-necessity result \cite{QG295} and the reconstruction audit
\cite{QG296}. Hence cosmology is consistent with the canonical architecture, completing the
physics part.
\end{proof}

\begin{corollary}[Difference → Actualization → Spectrum → Cosmology]
\label{cor:canonical-cosmology}
The cosmological sector completes the canonical chain:
\[
\text{Difference}\;\longrightarrow\;\text{Actualization}\;\longrightarrow\;\text{Inevitable Spectrum}\;\longrightarrow\;\text{Cosmology},
\]
with every cosmological observable a read of the spectrum.
\end{corollary}

\begin{proof}
By Theorem~\ref{thm:cosmo-consistent}, the cosmological observables are reads of the D96
spectrum; by Chapter~\ref{chap:spectrum} (Theorem~\ref{thm:core-completed}), the spectrum
is the output of the actualization process rooted in Difference. Hence the cosmological
sector is the terminal read of the canonical chain.
\end{proof}

\begin{remark}[Referee-safe wording]
Consistent with MONO007 \cite{MONO007}, this chapter presents cosmology as emergent and makes
no claims beyond the accepted derivations. No inflation claims beyond the accepted results
(the initial-condition and information origins replace the inflation epoch), no new
cosmological mechanisms, no unresolved Bekenstein discussion, and no boundary material are
included; the boundary topics are reserved for the Boundary Layer part of the monograph.
\end{remark}

\section{Summary}
\label{sec:cosmo-summary}

This chapter has derived cosmology. We have proved:

\begin{enumerate}
\item \textbf{Cosmology as a spectrum readout} (Theorem~\ref{thm:cosmo-emergent},
Corollary~\ref{cor:cosmo-emergent}): the cosmological observables emerge from the D96
spectral structure, and cosmology is emergent;
\item \textbf{The cosmological constant} (Theorem~\ref{thm:lambda},
Corollary~\ref{cor:lambda-emergent}): $\Lambda$ exists, is positive, and scales with the
actualization structure --- emergent, not fundamental;
\item \textbf{Matter and vacuum fractions} (Theorem~\ref{thm:fractions},
Corollary~\ref{cor:density-partition}): $\Omega_\Lambda$ and $\Omega_m$ are derived from
the D96 octave occupancies;
\item \textbf{The cosmic microwave background} (Theorem~\ref{thm:acoustic},
Corollary~\ref{cor:peak-hierarchy}, Lemma~\ref{lem:seed-spectrum}): the acoustic structure
is the D96 mode-ladder harmonics, with the seed spectrum the counting variance;
\item \textbf{Structure formation} (Theorem~\ref{thm:structure-formation},
Lemma~\ref{lem:clustering}): large-scale organization follows the actualization dynamics,
consistent with the spectral hierarchy;
\item \textbf{The expansion history} (Theorem~\ref{thm:expansion},
Corollary~\ref{cor:no-cosmo-primitive}): cosmological evolution is a spectrum consequence,
with no independent cosmological primitive;
\item \textbf{Consequences} (Theorem~\ref{thm:cosmo-consistent},
Corollary~\ref{cor:canonical-cosmology}): cosmology is consistent with the canonical
architecture, completing the physics part of the monograph.
\end{enumerate}

Cosmology is therefore emergent in The Actualization Theory: the cosmological constant is
the residual actualization pressure, the density fractions are the octave-record
information, the CMB acoustic structure is the spectral harmonics, structure formation
follows the actualization dynamics, and the expansion history is a spectrum consequence ---
with every observable traced back to the D96 spectral moments and no independent
cosmological primitive. This completes the physics part of the monograph
(Difference → Actualization → Spectrum → Physics): the foundation, the spectrum, quantum
mechanics, gravity, the standard model, and cosmology are all derived reads of the one
canonical structure.

\begin{thebibliography}{99}
\bibitem{QG227} AT-QG Phase 227, \emph{Initial Conditions Origin}, Docs/Research.
\bibitem{QG228} AT-QG Phase 228, \emph{Information Content Origin}, Docs/Research.
\bibitem{QG230} AT-QG Phase 230, \emph{Lambda Origin}, Docs/Research.
\bibitem{QG231} AT-QG Phase 231, \emph{Structure Formation Origin}, Docs/Research.
\bibitem{QG234} AT-QG Phase 234, \emph{Cosmological Fractions Origin}, Docs/Research.
\bibitem{QG237} AT-QG Phase 237, \emph{CMB Spectrum Origin}, Docs/Research.
\bibitem{QG238} AT-QG Phase 238, \emph{Acoustic Peak Origin}, Docs/Research.
\bibitem{QG240} AT-QG Phase 240, \emph{Cosmology Blind Reproduction}, Docs/Research.
\bibitem{QG295} AT-QG Phase 295, \emph{Spectrum Necessity Audit}, Docs/Research.
\bibitem{QG296} AT-QG Phase 296, \emph{Reconstruction Audit}, Docs/Research.
\bibitem{MONO007} AT-MONO007, \emph{Referee-Readiness Resolution}, Docs/Zenodo/Monograph.
\end{thebibliography}
